\section{Discussion: the object's proper name}
\label{sec:discussion}

Stepping back, the four registers of the paper describe one thing, and that thing has a
proper name in mathematics. It is a \emph{metallic-mean Schr\"odinger cocycle} over the
substitution subshift generated by $R^mL^m$, analysed by the Kohmoto--Kadanoff--Tang
trace map, with $\kk$ its Fricke--Vogt invariant \cite{kkt1983,suto1987,baakegrimm2013}.

\subsection{Non-cancellation as positivity}

The founding intuition --- that ``not-nothing'' should be a residue that cannot cancel to
zero --- has an exact translation in this language. Non-cancellation is Fricke--Vogt
positivity $\kk>2$; this is precisely the hyperbolicity condition of the trace map; and
that hyperbolicity is exactly the condition for the associated one-dimensional aperiodic
Schr\"odinger operator to have purely singular-continuous spectrum supported on a Cantor
set \cite{suto1987}. Figure~\ref{fig:cantor} shows that spectrum for the golden case: as
the coupling grows, gaps open on every scale and the spectrum becomes a Cantor set --- the
operator-theoretic shadow of the metallic dynamics. This is genuine, observed physics ---
aperiodic order, the spectral theory of quasicrystals --- and it is emphatically
\emph{not} fundamental physics: there is no gauge theory here, no constant of nature, no
spacetime. It is the honest physical home of the object, and it sits comfortably on the
mathematics side of the firewall.

\begin{figure}[t]
  \centering
  \includegraphics[width=0.86\linewidth]{fig11_cantor_spectrum}
  \caption{The object's proper name. The spectrum of the Fibonacci (golden-mean)
  Schr\"odinger operator as a function of the coupling $V$: a Cantor set whose gaps are
  organised by the trace-map dynamics. Non-cancellation $\kk>2$ is exactly the
  hyperbolicity that produces it. \sNum}
  \label{fig:cantor}
\end{figure}

\subsection{The obstruction map}

Reading the paper as a whole, the object is bounded by a small number of clearly located
walls, and naming them is as much a result as the bridges. The motivating ambition met
obstructions at six recurring places: a \emph{source} (what generates the family is a
chosen minimality, not a derived necessity); a \emph{selector} (no principle forces the
golden seed over the others); a \emph{measure} (the character variety parametrises
possibilities, it does not weight them); a \emph{unit} (the invariants are dimensionless,
Section~\ref{sec:bridge}); a \emph{dictionary} (the class-$S$ bridge is a symmetry
identity, not a field-content map); and an \emph{observable} (nothing in the construction
inserts a measurement). The firewall is the discipline of mapping these walls honestly,
and the bridge/wall pairing of Section~\ref{sec:bridge} is the sharpest instance: the
object reaches as far as a symmetry and stops exactly at a scale.

\subsection{What the object is, and is not}

It is a small, self-generating mathematical object that turned out to know about knots
(the figure-eight $A$-polynomial), Lie theory (the Dickson tower and the opposition
involution), aperiodic order (the Cantor spectrum), and a gauge theory's duality (the
class-$S$ bridge). It is not a theory of physical values; the firewall holds, and we have
exhibited why. The most honest summary is that the object knows, precisely, where its own
reach ends --- and that this boundary, mapped exactly, is the paper's real subject.
